\chapter{Operaciones CRUD con JPA/ORM validadas con Postman}
\label{ch:crud}

\section{Persistencia con Spring Data JPA}

Cada microservicio modela su dominio como una entidad JPA independiente, persistida en
su propia base de datos mediante \texttt{JpaRepository}, sin necesidad de escribir SQL
manual para las operaciones CRUD básicas. El siguiente ejemplo corresponde a la entidad
\texttt{Libro} del microservicio de catálogo:

\begin{lstlisting}[language=Java, caption={Libro.java -- entidad JPA de catálogo}]
@Entity
@Table(name = "libros")
@Data
public class Libro {
    @Id
    @GeneratedValue(strategy = GenerationType.IDENTITY)
    private Long id;

    private String titulo;
    private String autor;
    private Double precio;
}
\end{lstlisting}

\begin{lstlisting}[language=Java, caption={LibroRepository.java -- repositorio}]
public interface LibroRepository extends JpaRepository<Libro, Long> {
}
\end{lstlisting}

La sola extensión de \texttt{JpaRepository<Libro, Long>} otorga las operaciones
\texttt{save}, \texttt{findById}, \texttt{findAll} y \texttt{deleteById} sin
implementación adicional. En los servicios donde la consulta lo requiere se agregan
métodos derivados (por ejemplo, \texttt{findByLibroIdAndSucursal} en inventario, o
\texttt{findByCodigoSeguimiento} en envíos), que Spring Data traduce automáticamente a
JPQL a partir del nombre del método.

Sobre esta base, cada servicio expone las cuatro operaciones CRUD que correspondan a su
dominio a través de endpoints REST. Por ejemplo, el microservicio de clientes expone
las cuatro operaciones completas:

\begin{center}
\begin{tabular}{lll}
\toprule
\textbf{Operación} & \textbf{Método HTTP} & \textbf{Endpoint} \\
\midrule
Crear & POST & \texttt{/api/v1/clientes} \\
Leer (uno) & GET & \texttt{/api/v1/clientes/\{id\}} \\
Leer (todos) & GET & \texttt{/api/v1/clientes} \\
Actualizar & PUT & \texttt{/api/v1/clientes/\{id\}} \\
Eliminar & DELETE & \texttt{/api/v1/clientes/\{id\}} \\
\bottomrule
\end{tabular}
\end{center}

Otros servicios exponen un subconjunto de estas operaciones cuando el dominio no
requiere todas -- por ejemplo, catálogo no expone \texttt{DELETE} porque un libro
referenciado por pedidos o reseñas no debe poder eliminarse sin más.

\section{Validación con la colección Postman}

Para cumplir con el requisito de validar las operaciones CRUD con datos provenientes de
Postman, se construyó una colección (\texttt{postman/BookPoint.postman\_collection.json})
que cubre los diez microservicios, documentada en \texttt{postman/README.md}.

La colección está organizada en once carpetas ordenadas según la cadena real de
dependencias entre servicios (catálogo $\rightarrow$ clientes $\rightarrow$ inventario
$\rightarrow$ carrito $\rightarrow$ pedidos $\rightarrow$ pagos $\rightarrow$ envíos
$\rightarrow$ reseñas $\rightarrow$ notificaciones $\rightarrow$ sucursales), de modo
que puede ejecutarse de principio a fin con \textit{Collection Runner}. Cada request de
creación incluye un script de prueba que captura el \texttt{id} devuelto por la API y
lo guarda en una variable de colección, reutilizada automáticamente por los siguientes
requests -- por ejemplo, el \texttt{id} del libro creado en la carpeta de catálogo se
reutiliza al registrar stock en inventario, agregarlo al carrito y generar un pedido.

\begin{lstlisting}[caption={Script de test de Postman -- captura automática de ID al crear un libro}]
pm.test('Status 201', () => pm.response.to.have.status(201));
const json = pm.response.json();
pm.collectionVariables.set('libroId', json.id);
\end{lstlisting}

Todas las requests de la colección fueron probadas en vivo contra los diez servicios
levantados en Docker antes de considerarlas válidas -- no se trata de una colección
escrita en teoría, sino verificada extremo a extremo, incluyendo los datos generados
automáticamente por Datafaker como insumo (sección~\ref{sec:evolucion}).

\section{Reportes de ventas}

El caso pide que el perfil Jefe de Sucursal pueda "generar reportes de ventas por
categoría, autor o sucursal". Se implementaron los dos primeros en el microservicio de
pedidos, que ya concentra los datos de venta (\texttt{sucursal}, \texttt{libroId},
\texttt{total}):

\begin{itemize}
    \item \texttt{GET /api/v1/pedidos/reportes/por-sucursal}: agregación directa sobre
    la propia base de datos de pedidos (\texttt{GROUP BY sucursal} vía una consulta
    JPQL con proyección de interfaz de Spring Data).
    \item \texttt{GET /api/v1/pedidos/reportes/por-autor}: como el autor no vive en la
    base de datos de pedidos sino en la de catálogo, no es posible resolverlo con un
    \texttt{JOIN} SQL. Se agregan primero las ventas por \texttt{libroId}, y luego se
    enriquece cada grupo con el autor correspondiente vía Feign, reagrupando en memoria
    los libros que comparten autor.
\end{itemize}

El reporte por categoría se dejó fuera: exigiría agregar un campo de categoría o
género al catálogo, un cambio de esquema en otro microservicio no motivado por ninguna
otra necesidad actual del sistema.

Se verificó en vivo, contra el stack completo levantado en Docker con datos generados
por Datafaker, que ambos reportes son consistentes entre sí: la suma total de ventas
agrupada por sucursal coincide exactamente con la suma total agrupada por autor
(\$796.892 en ambos casos sobre el mismo conjunto de pedidos), confirmando que ninguna
venta se pierde ni se duplica al reagrupar.
